\documentclass[11pt,a4paper]{article}
% main (standalone preamble for editing)
% vim: nowrap: spell:
% ══════════════════════════════════════════════════════════════════════════════
% Speed Maths --- shared preamble
% Included by every sheet and answer file via:  % main (standalone preamble for editing)
% vim: nowrap: spell:
% ══════════════════════════════════════════════════════════════════════════════
% Speed Maths --- shared preamble
% Included by every sheet and answer file via:  % main (standalone preamble for editing)
% vim: nowrap: spell:
% ══════════════════════════════════════════════════════════════════════════════
% Speed Maths --- shared preamble
% Included by every sheet and answer file via:  \input{../../shared/preamble}
% Editing THIS file changes the look of every sheet/answer across every pillar.
% ══════════════════════════════════════════════════════════════════════════════

\usepackage[a4paper, top=25mm, bottom=25mm, left=22mm, right=22mm]{geometry}
\usepackage{amsmath, amssymb}
\usepackage{enumitem}
\usepackage{booktabs}
\usepackage{array}
\usepackage{parskip}
\usepackage{microtype}
\usepackage{fancyhdr}
\usepackage{titlesec}
\usepackage{mdframed}
\usepackage{xcolor}
\usepackage[hidelinks]{hyperref}
\usepackage{graphicx}
\usepackage{tikz}
\usepackage{pgfplots}
\pgfplotsset{compat=1.18}

% ── Colours ───────────────────────────────────────────────────────────────────
\definecolor{darkgray}{gray}{0.15}
\definecolor{medgray}{gray}{0.45}
\definecolor{lightgray}{gray}{0.93}
\definecolor{boxgray}{gray}{0.88}

% ── Section formatting ──────────────────────────────────────────────────────────
\titleformat{\section}[block]
  {\normalfont\large\bfseries}{}
  {0pt}{}[\vspace{2pt}\hrule\vspace{6pt}]
\titlespacing*{\section}{0pt}{14pt}{4pt}

% ── List formatting ─────────────────────────────────────────────────────────────
\setlist[enumerate,1]{label=\textbf{\arabic*.}, leftmargin=2em, itemsep=10pt}

% ── Branding constants (edit here, not per-file) ────────────────────────────────
\newcommand{\SpeedExamLine}{\textbf{Speed Maths:} Competition \& Exam Prep}
\newcommand{\SpeedCredit}{Series by \href{https://www.linkedin.com/in/cerealdev/}{David Akinyele-Aje}}

% ── Header / footer ──────────────────────────────────────────────────────────────
% Usage (once, after \input{../../shared/preamble} and before \begin{document}):
%   \SpeedHeader{Algebra}{3}
% First arg  = pillar name shown as "Speed <Pillar> --- Daily Drill"
% Second arg = worksheet number
\pagestyle{fancy}
\newcommand{\SpeedHeader}[2]{%
  \fancyhf{}%
  \renewcommand{\headrulewidth}{0.4pt}%
  \setlength{\headheight}{14pt}%
  \fancyhead[L]{\small\textbf{Speed #1 --- Daily Drill}}%
  \fancyhead[R]{\small Worksheet \##2}%
  \fancyfoot[L]{\small \SpeedExamLine}%
  \fancyfoot[C]{\small\thepage}%
  \fancyfoot[R]{\small \SpeedCredit}%
  \renewcommand{\footrulewidth}{0.4pt}%
}

% ── Title block (used at the top of every sheet AND answer file) ───────────────
% Usage: \SpeedTitleBlock{Daily Algebra Drill \#3}{Jane Doe}
% %% AGENTS: When generating a new sheet, ask the human contributor for the name they want credited in argument 2.
% For answer files, pass the "--- Answers \& Investigations" suffix in the arg.
\newcommand{\SpeedTitleBlock}[2]{%
  \begin{center}
    {\LARGE\bfseries #1}\\[4pt]
    {\normalsize\itshape Sheet by #2}\\[6pt]
    \hrule\vspace{4pt}
    \begin{tabular}{llcc}
      \toprule
      \textbf{Section} & \textbf{Topic} & \textbf{Questions} & \textbf{Target time}\\
      \midrule
      A & Rapid Recognition          & 10 & 2 min 30 sec \\
      B & Manipulation Drills        & 10 & 8 min        \\
      C & Substitution \& Structure  & 8  & 10 min       \\
      D & Challenge                  & 5  & 15 min       \\
      \bottomrule
    \end{tabular}
    \vspace{4pt}
    \hrule
  \end{center}
}

% ── Sheet metadata: topic + the day's new tools ────────────────────────────────
% Usage (immediately after \SpeedTitleBlock, sheets only):
%   \SpeedMeta{Expanding, factorising, and the identity toolkit}
%             {difference of squares, sum/product form, completing the square}
%
% Arg 1 = the sheet's topic in one line. The website lists this instead of
%         "Sheet 03", so write the thing a reader would actually search for.
% Arg 2 = comma-separated, at most five: the tools this sheet introduces that
%         earlier sheets in the pillar did not. These become the website's tags.
%         Leave out anything the reader already met — the tags are meant to say
%         what is new today, not to summarise the sheet.
%
% Both are parsed straight out of this file by tools/build_website.py, so the
% sheet stays the single source of truth and the site cannot drift from it.
%
% This renders NOTHING. It is a declaration for the build script, not a piece of
% the sheet's layout: adding it to a sheet must not change that sheet's PDF, so
% the 25 existing sheets can be annotated without a single one of them being
% reissued. If the toolkit should also appear on the page, that is a separate
% decision about the sheet design — make it deliberately, here, not by leaving
% this macro printing something nobody asked it to.
\newcommand{\SpeedMeta}[2]{}

% ── Answer / Method / Investigate-further boxes (answer files only) ────────────
\newmdenv[
  backgroundcolor=lightgray,
  linecolor=medgray,
  linewidth=0.5pt,
  leftmargin=0pt, rightmargin=0pt,
  innerleftmargin=8pt, innerrightmargin=8pt,
  innertopmargin=5pt, innerbottommargin=5pt,
  skipabove=4pt, skipbelow=4pt
]{ansbox}

\newmdenv[
  backgroundcolor=white,
  linecolor=darkgray,
  linewidth=0.8pt,
  leftline=true, rightline=false, topline=false, bottomline=false,
  leftmargin=0pt, rightmargin=0pt,
  innerleftmargin=8pt, innerrightmargin=6pt,
  innertopmargin=4pt, innerbottommargin=4pt,
  skipabove=4pt, skipbelow=6pt
]{invbox}

\newcommand{\ans}[1]{%
  \begin{ansbox}
    \textbf{Answer:} #1
  \end{ansbox}}

\newcommand{\method}[1]{%
  {\small\textbf{Method:} #1}\par}

\newcommand{\inv}[1]{%
  \begin{invbox}
    \textbf{\small Investigate further:} {\small #1}
  \end{invbox}}

% ── Misc helpers ────────────────────────────────────────────────────────────────
\newcommand{\qn}[1]{\textbf{#1.}\;}

% Mistake-linking: attach a Top 5 Patterns / Common Traps bullet to the question
% label(s) in this sheet where it actually shows up, e.g. \seealso{B6, D2}
\newcommand{\seealso}[1]{\hfill\textit{\footnotesize(see #1)}}

% ── Graph options macro ─────────────────────────────────────────────────────────
% Usage: \GraphOptions{tikz code for A}{tikz code for B}{tikz code for C}{tikz code for D}
\newcommand{\GraphOptions}[4]{%
  \begin{center}
    \begin{tabular}{cc}
      \textbf{A)} & \textbf{B)} \\
      #1 & #2 \\[10pt]
      \textbf{C)} & \textbf{D)} \\
      #3 & #4
    \end{tabular}
  \end{center}
}

% ── Closing motivational rule + cross-promo footer (sheets only) ────────────────
% Usage: \SpeedClosing{"Quote text."}
\newcommand{\SpeedClosing}[1]{%
  \vspace{20pt}
  \begin{center}
  \rule{0.6\textwidth}{0.4pt}\\[4pt]
  {\small\itshape #1}\\[4pt]
  \rule{0.6\textwidth}{0.4pt}
  \end{center}
  \begin{center}
  {\small Found a mistake or want to contribute a sheet?}\\
  {\small Visit the open-source project at \href{https://github.com/The-CerealDev/speed-maths}{github.com/The-CerealDev/speed-maths}}
  \end{center}
}

% Editing THIS file changes the look of every sheet/answer across every pillar.
% ══════════════════════════════════════════════════════════════════════════════

\usepackage[a4paper, top=25mm, bottom=25mm, left=22mm, right=22mm]{geometry}
\usepackage{amsmath, amssymb}
\usepackage{enumitem}
\usepackage{booktabs}
\usepackage{array}
\usepackage{parskip}
\usepackage{microtype}
\usepackage{fancyhdr}
\usepackage{titlesec}
\usepackage{mdframed}
\usepackage{xcolor}
\usepackage[hidelinks]{hyperref}
\usepackage{graphicx}
\usepackage{tikz}
\usepackage{pgfplots}
\pgfplotsset{compat=1.18}

% ── Colours ───────────────────────────────────────────────────────────────────
\definecolor{darkgray}{gray}{0.15}
\definecolor{medgray}{gray}{0.45}
\definecolor{lightgray}{gray}{0.93}
\definecolor{boxgray}{gray}{0.88}

% ── Section formatting ──────────────────────────────────────────────────────────
\titleformat{\section}[block]
  {\normalfont\large\bfseries}{}
  {0pt}{}[\vspace{2pt}\hrule\vspace{6pt}]
\titlespacing*{\section}{0pt}{14pt}{4pt}

% ── List formatting ─────────────────────────────────────────────────────────────
\setlist[enumerate,1]{label=\textbf{\arabic*.}, leftmargin=2em, itemsep=10pt}

% ── Branding constants (edit here, not per-file) ────────────────────────────────
\newcommand{\SpeedExamLine}{\textbf{Speed Maths:} Competition \& Exam Prep}
\newcommand{\SpeedCredit}{Series by \href{https://www.linkedin.com/in/cerealdev/}{David Akinyele-Aje}}

% ── Header / footer ──────────────────────────────────────────────────────────────
% Usage (once, after % main (standalone preamble for editing)
% vim: nowrap: spell:
% ══════════════════════════════════════════════════════════════════════════════
% Speed Maths --- shared preamble
% Included by every sheet and answer file via:  \input{../../shared/preamble}
% Editing THIS file changes the look of every sheet/answer across every pillar.
% ══════════════════════════════════════════════════════════════════════════════

\usepackage[a4paper, top=25mm, bottom=25mm, left=22mm, right=22mm]{geometry}
\usepackage{amsmath, amssymb}
\usepackage{enumitem}
\usepackage{booktabs}
\usepackage{array}
\usepackage{parskip}
\usepackage{microtype}
\usepackage{fancyhdr}
\usepackage{titlesec}
\usepackage{mdframed}
\usepackage{xcolor}
\usepackage[hidelinks]{hyperref}
\usepackage{graphicx}
\usepackage{tikz}
\usepackage{pgfplots}
\pgfplotsset{compat=1.18}

% ── Colours ───────────────────────────────────────────────────────────────────
\definecolor{darkgray}{gray}{0.15}
\definecolor{medgray}{gray}{0.45}
\definecolor{lightgray}{gray}{0.93}
\definecolor{boxgray}{gray}{0.88}

% ── Section formatting ──────────────────────────────────────────────────────────
\titleformat{\section}[block]
  {\normalfont\large\bfseries}{}
  {0pt}{}[\vspace{2pt}\hrule\vspace{6pt}]
\titlespacing*{\section}{0pt}{14pt}{4pt}

% ── List formatting ─────────────────────────────────────────────────────────────
\setlist[enumerate,1]{label=\textbf{\arabic*.}, leftmargin=2em, itemsep=10pt}

% ── Branding constants (edit here, not per-file) ────────────────────────────────
\newcommand{\SpeedExamLine}{\textbf{Speed Maths:} Competition \& Exam Prep}
\newcommand{\SpeedCredit}{Series by \href{https://www.linkedin.com/in/cerealdev/}{David Akinyele-Aje}}

% ── Header / footer ──────────────────────────────────────────────────────────────
% Usage (once, after \input{../../shared/preamble} and before \begin{document}):
%   \SpeedHeader{Algebra}{3}
% First arg  = pillar name shown as "Speed <Pillar> --- Daily Drill"
% Second arg = worksheet number
\pagestyle{fancy}
\newcommand{\SpeedHeader}[2]{%
  \fancyhf{}%
  \renewcommand{\headrulewidth}{0.4pt}%
  \setlength{\headheight}{14pt}%
  \fancyhead[L]{\small\textbf{Speed #1 --- Daily Drill}}%
  \fancyhead[R]{\small Worksheet \##2}%
  \fancyfoot[L]{\small \SpeedExamLine}%
  \fancyfoot[C]{\small\thepage}%
  \fancyfoot[R]{\small \SpeedCredit}%
  \renewcommand{\footrulewidth}{0.4pt}%
}

% ── Title block (used at the top of every sheet AND answer file) ───────────────
% Usage: \SpeedTitleBlock{Daily Algebra Drill \#3}{Jane Doe}
% %% AGENTS: When generating a new sheet, ask the human contributor for the name they want credited in argument 2.
% For answer files, pass the "--- Answers \& Investigations" suffix in the arg.
\newcommand{\SpeedTitleBlock}[2]{%
  \begin{center}
    {\LARGE\bfseries #1}\\[4pt]
    {\normalsize\itshape Sheet by #2}\\[6pt]
    \hrule\vspace{4pt}
    \begin{tabular}{llcc}
      \toprule
      \textbf{Section} & \textbf{Topic} & \textbf{Questions} & \textbf{Target time}\\
      \midrule
      A & Rapid Recognition          & 10 & 2 min 30 sec \\
      B & Manipulation Drills        & 10 & 8 min        \\
      C & Substitution \& Structure  & 8  & 10 min       \\
      D & Challenge                  & 5  & 15 min       \\
      \bottomrule
    \end{tabular}
    \vspace{4pt}
    \hrule
  \end{center}
}

% ── Sheet metadata: topic + the day's new tools ────────────────────────────────
% Usage (immediately after \SpeedTitleBlock, sheets only):
%   \SpeedMeta{Expanding, factorising, and the identity toolkit}
%             {difference of squares, sum/product form, completing the square}
%
% Arg 1 = the sheet's topic in one line. The website lists this instead of
%         "Sheet 03", so write the thing a reader would actually search for.
% Arg 2 = comma-separated, at most five: the tools this sheet introduces that
%         earlier sheets in the pillar did not. These become the website's tags.
%         Leave out anything the reader already met — the tags are meant to say
%         what is new today, not to summarise the sheet.
%
% Both are parsed straight out of this file by tools/build_website.py, so the
% sheet stays the single source of truth and the site cannot drift from it.
%
% This renders NOTHING. It is a declaration for the build script, not a piece of
% the sheet's layout: adding it to a sheet must not change that sheet's PDF, so
% the 25 existing sheets can be annotated without a single one of them being
% reissued. If the toolkit should also appear on the page, that is a separate
% decision about the sheet design — make it deliberately, here, not by leaving
% this macro printing something nobody asked it to.
\newcommand{\SpeedMeta}[2]{}

% ── Answer / Method / Investigate-further boxes (answer files only) ────────────
\newmdenv[
  backgroundcolor=lightgray,
  linecolor=medgray,
  linewidth=0.5pt,
  leftmargin=0pt, rightmargin=0pt,
  innerleftmargin=8pt, innerrightmargin=8pt,
  innertopmargin=5pt, innerbottommargin=5pt,
  skipabove=4pt, skipbelow=4pt
]{ansbox}

\newmdenv[
  backgroundcolor=white,
  linecolor=darkgray,
  linewidth=0.8pt,
  leftline=true, rightline=false, topline=false, bottomline=false,
  leftmargin=0pt, rightmargin=0pt,
  innerleftmargin=8pt, innerrightmargin=6pt,
  innertopmargin=4pt, innerbottommargin=4pt,
  skipabove=4pt, skipbelow=6pt
]{invbox}

\newcommand{\ans}[1]{%
  \begin{ansbox}
    \textbf{Answer:} #1
  \end{ansbox}}

\newcommand{\method}[1]{%
  {\small\textbf{Method:} #1}\par}

\newcommand{\inv}[1]{%
  \begin{invbox}
    \textbf{\small Investigate further:} {\small #1}
  \end{invbox}}

% ── Misc helpers ────────────────────────────────────────────────────────────────
\newcommand{\qn}[1]{\textbf{#1.}\;}

% Mistake-linking: attach a Top 5 Patterns / Common Traps bullet to the question
% label(s) in this sheet where it actually shows up, e.g. \seealso{B6, D2}
\newcommand{\seealso}[1]{\hfill\textit{\footnotesize(see #1)}}

% ── Graph options macro ─────────────────────────────────────────────────────────
% Usage: \GraphOptions{tikz code for A}{tikz code for B}{tikz code for C}{tikz code for D}
\newcommand{\GraphOptions}[4]{%
  \begin{center}
    \begin{tabular}{cc}
      \textbf{A)} & \textbf{B)} \\
      #1 & #2 \\[10pt]
      \textbf{C)} & \textbf{D)} \\
      #3 & #4
    \end{tabular}
  \end{center}
}

% ── Closing motivational rule + cross-promo footer (sheets only) ────────────────
% Usage: \SpeedClosing{"Quote text."}
\newcommand{\SpeedClosing}[1]{%
  \vspace{20pt}
  \begin{center}
  \rule{0.6\textwidth}{0.4pt}\\[4pt]
  {\small\itshape #1}\\[4pt]
  \rule{0.6\textwidth}{0.4pt}
  \end{center}
  \begin{center}
  {\small Found a mistake or want to contribute a sheet?}\\
  {\small Visit the open-source project at \href{https://github.com/The-CerealDev/speed-maths}{github.com/The-CerealDev/speed-maths}}
  \end{center}
}
 and before \begin{document}):
%   \SpeedHeader{Algebra}{3}
% First arg  = pillar name shown as "Speed <Pillar> --- Daily Drill"
% Second arg = worksheet number
\pagestyle{fancy}
\newcommand{\SpeedHeader}[2]{%
  \fancyhf{}%
  \renewcommand{\headrulewidth}{0.4pt}%
  \setlength{\headheight}{14pt}%
  \fancyhead[L]{\small\textbf{Speed #1 --- Daily Drill}}%
  \fancyhead[R]{\small Worksheet \##2}%
  \fancyfoot[L]{\small \SpeedExamLine}%
  \fancyfoot[C]{\small\thepage}%
  \fancyfoot[R]{\small \SpeedCredit}%
  \renewcommand{\footrulewidth}{0.4pt}%
}

% ── Title block (used at the top of every sheet AND answer file) ───────────────
% Usage: \SpeedTitleBlock{Daily Algebra Drill \#3}{Jane Doe}
% %% AGENTS: When generating a new sheet, ask the human contributor for the name they want credited in argument 2.
% For answer files, pass the "--- Answers \& Investigations" suffix in the arg.
\newcommand{\SpeedTitleBlock}[2]{%
  \begin{center}
    {\LARGE\bfseries #1}\\[4pt]
    {\normalsize\itshape Sheet by #2}\\[6pt]
    \hrule\vspace{4pt}
    \begin{tabular}{llcc}
      \toprule
      \textbf{Section} & \textbf{Topic} & \textbf{Questions} & \textbf{Target time}\\
      \midrule
      A & Rapid Recognition          & 10 & 2 min 30 sec \\
      B & Manipulation Drills        & 10 & 8 min        \\
      C & Substitution \& Structure  & 8  & 10 min       \\
      D & Challenge                  & 5  & 15 min       \\
      \bottomrule
    \end{tabular}
    \vspace{4pt}
    \hrule
  \end{center}
}

% ── Sheet metadata: topic + the day's new tools ────────────────────────────────
% Usage (immediately after \SpeedTitleBlock, sheets only):
%   \SpeedMeta{Expanding, factorising, and the identity toolkit}
%             {difference of squares, sum/product form, completing the square}
%
% Arg 1 = the sheet's topic in one line. The website lists this instead of
%         "Sheet 03", so write the thing a reader would actually search for.
% Arg 2 = comma-separated, at most five: the tools this sheet introduces that
%         earlier sheets in the pillar did not. These become the website's tags.
%         Leave out anything the reader already met — the tags are meant to say
%         what is new today, not to summarise the sheet.
%
% Both are parsed straight out of this file by tools/build_website.py, so the
% sheet stays the single source of truth and the site cannot drift from it.
%
% This renders NOTHING. It is a declaration for the build script, not a piece of
% the sheet's layout: adding it to a sheet must not change that sheet's PDF, so
% the 25 existing sheets can be annotated without a single one of them being
% reissued. If the toolkit should also appear on the page, that is a separate
% decision about the sheet design — make it deliberately, here, not by leaving
% this macro printing something nobody asked it to.
\newcommand{\SpeedMeta}[2]{}

% ── Answer / Method / Investigate-further boxes (answer files only) ────────────
\newmdenv[
  backgroundcolor=lightgray,
  linecolor=medgray,
  linewidth=0.5pt,
  leftmargin=0pt, rightmargin=0pt,
  innerleftmargin=8pt, innerrightmargin=8pt,
  innertopmargin=5pt, innerbottommargin=5pt,
  skipabove=4pt, skipbelow=4pt
]{ansbox}

\newmdenv[
  backgroundcolor=white,
  linecolor=darkgray,
  linewidth=0.8pt,
  leftline=true, rightline=false, topline=false, bottomline=false,
  leftmargin=0pt, rightmargin=0pt,
  innerleftmargin=8pt, innerrightmargin=6pt,
  innertopmargin=4pt, innerbottommargin=4pt,
  skipabove=4pt, skipbelow=6pt
]{invbox}

\newcommand{\ans}[1]{%
  \begin{ansbox}
    \textbf{Answer:} #1
  \end{ansbox}}

\newcommand{\method}[1]{%
  {\small\textbf{Method:} #1}\par}

\newcommand{\inv}[1]{%
  \begin{invbox}
    \textbf{\small Investigate further:} {\small #1}
  \end{invbox}}

% ── Misc helpers ────────────────────────────────────────────────────────────────
\newcommand{\qn}[1]{\textbf{#1.}\;}

% Mistake-linking: attach a Top 5 Patterns / Common Traps bullet to the question
% label(s) in this sheet where it actually shows up, e.g. \seealso{B6, D2}
\newcommand{\seealso}[1]{\hfill\textit{\footnotesize(see #1)}}

% ── Graph options macro ─────────────────────────────────────────────────────────
% Usage: \GraphOptions{tikz code for A}{tikz code for B}{tikz code for C}{tikz code for D}
\newcommand{\GraphOptions}[4]{%
  \begin{center}
    \begin{tabular}{cc}
      \textbf{A)} & \textbf{B)} \\
      #1 & #2 \\[10pt]
      \textbf{C)} & \textbf{D)} \\
      #3 & #4
    \end{tabular}
  \end{center}
}

% ── Closing motivational rule + cross-promo footer (sheets only) ────────────────
% Usage: \SpeedClosing{"Quote text."}
\newcommand{\SpeedClosing}[1]{%
  \vspace{20pt}
  \begin{center}
  \rule{0.6\textwidth}{0.4pt}\\[4pt]
  {\small\itshape #1}\\[4pt]
  \rule{0.6\textwidth}{0.4pt}
  \end{center}
  \begin{center}
  {\small Found a mistake or want to contribute a sheet?}\\
  {\small Visit the open-source project at \href{https://github.com/The-CerealDev/speed-maths}{github.com/The-CerealDev/speed-maths}}
  \end{center}
}

% Editing THIS file changes the look of every sheet/answer across every pillar.
% ══════════════════════════════════════════════════════════════════════════════

\usepackage[a4paper, top=25mm, bottom=25mm, left=22mm, right=22mm]{geometry}
\usepackage{amsmath, amssymb}
\usepackage{enumitem}
\usepackage{booktabs}
\usepackage{array}
\usepackage{parskip}
\usepackage{microtype}
\usepackage{fancyhdr}
\usepackage{titlesec}
\usepackage{mdframed}
\usepackage{xcolor}
\usepackage[hidelinks]{hyperref}
\usepackage{graphicx}
\usepackage{tikz}
\usepackage{pgfplots}
\pgfplotsset{compat=1.18}

% ── Colours ───────────────────────────────────────────────────────────────────
\definecolor{darkgray}{gray}{0.15}
\definecolor{medgray}{gray}{0.45}
\definecolor{lightgray}{gray}{0.93}
\definecolor{boxgray}{gray}{0.88}

% ── Section formatting ──────────────────────────────────────────────────────────
\titleformat{\section}[block]
  {\normalfont\large\bfseries}{}
  {0pt}{}[\vspace{2pt}\hrule\vspace{6pt}]
\titlespacing*{\section}{0pt}{14pt}{4pt}

% ── List formatting ─────────────────────────────────────────────────────────────
\setlist[enumerate,1]{label=\textbf{\arabic*.}, leftmargin=2em, itemsep=10pt}

% ── Branding constants (edit here, not per-file) ────────────────────────────────
\newcommand{\SpeedExamLine}{\textbf{Speed Maths:} Competition \& Exam Prep}
\newcommand{\SpeedCredit}{Series by \href{https://www.linkedin.com/in/cerealdev/}{David Akinyele-Aje}}

% ── Header / footer ──────────────────────────────────────────────────────────────
% Usage (once, after % main (standalone preamble for editing)
% vim: nowrap: spell:
% ══════════════════════════════════════════════════════════════════════════════
% Speed Maths --- shared preamble
% Included by every sheet and answer file via:  % main (standalone preamble for editing)
% vim: nowrap: spell:
% ══════════════════════════════════════════════════════════════════════════════
% Speed Maths --- shared preamble
% Included by every sheet and answer file via:  \input{../../shared/preamble}
% Editing THIS file changes the look of every sheet/answer across every pillar.
% ══════════════════════════════════════════════════════════════════════════════

\usepackage[a4paper, top=25mm, bottom=25mm, left=22mm, right=22mm]{geometry}
\usepackage{amsmath, amssymb}
\usepackage{enumitem}
\usepackage{booktabs}
\usepackage{array}
\usepackage{parskip}
\usepackage{microtype}
\usepackage{fancyhdr}
\usepackage{titlesec}
\usepackage{mdframed}
\usepackage{xcolor}
\usepackage[hidelinks]{hyperref}
\usepackage{graphicx}
\usepackage{tikz}
\usepackage{pgfplots}
\pgfplotsset{compat=1.18}

% ── Colours ───────────────────────────────────────────────────────────────────
\definecolor{darkgray}{gray}{0.15}
\definecolor{medgray}{gray}{0.45}
\definecolor{lightgray}{gray}{0.93}
\definecolor{boxgray}{gray}{0.88}

% ── Section formatting ──────────────────────────────────────────────────────────
\titleformat{\section}[block]
  {\normalfont\large\bfseries}{}
  {0pt}{}[\vspace{2pt}\hrule\vspace{6pt}]
\titlespacing*{\section}{0pt}{14pt}{4pt}

% ── List formatting ─────────────────────────────────────────────────────────────
\setlist[enumerate,1]{label=\textbf{\arabic*.}, leftmargin=2em, itemsep=10pt}

% ── Branding constants (edit here, not per-file) ────────────────────────────────
\newcommand{\SpeedExamLine}{\textbf{Speed Maths:} Competition \& Exam Prep}
\newcommand{\SpeedCredit}{Series by \href{https://www.linkedin.com/in/cerealdev/}{David Akinyele-Aje}}

% ── Header / footer ──────────────────────────────────────────────────────────────
% Usage (once, after \input{../../shared/preamble} and before \begin{document}):
%   \SpeedHeader{Algebra}{3}
% First arg  = pillar name shown as "Speed <Pillar> --- Daily Drill"
% Second arg = worksheet number
\pagestyle{fancy}
\newcommand{\SpeedHeader}[2]{%
  \fancyhf{}%
  \renewcommand{\headrulewidth}{0.4pt}%
  \setlength{\headheight}{14pt}%
  \fancyhead[L]{\small\textbf{Speed #1 --- Daily Drill}}%
  \fancyhead[R]{\small Worksheet \##2}%
  \fancyfoot[L]{\small \SpeedExamLine}%
  \fancyfoot[C]{\small\thepage}%
  \fancyfoot[R]{\small \SpeedCredit}%
  \renewcommand{\footrulewidth}{0.4pt}%
}

% ── Title block (used at the top of every sheet AND answer file) ───────────────
% Usage: \SpeedTitleBlock{Daily Algebra Drill \#3}{Jane Doe}
% %% AGENTS: When generating a new sheet, ask the human contributor for the name they want credited in argument 2.
% For answer files, pass the "--- Answers \& Investigations" suffix in the arg.
\newcommand{\SpeedTitleBlock}[2]{%
  \begin{center}
    {\LARGE\bfseries #1}\\[4pt]
    {\normalsize\itshape Sheet by #2}\\[6pt]
    \hrule\vspace{4pt}
    \begin{tabular}{llcc}
      \toprule
      \textbf{Section} & \textbf{Topic} & \textbf{Questions} & \textbf{Target time}\\
      \midrule
      A & Rapid Recognition          & 10 & 2 min 30 sec \\
      B & Manipulation Drills        & 10 & 8 min        \\
      C & Substitution \& Structure  & 8  & 10 min       \\
      D & Challenge                  & 5  & 15 min       \\
      \bottomrule
    \end{tabular}
    \vspace{4pt}
    \hrule
  \end{center}
}

% ── Sheet metadata: topic + the day's new tools ────────────────────────────────
% Usage (immediately after \SpeedTitleBlock, sheets only):
%   \SpeedMeta{Expanding, factorising, and the identity toolkit}
%             {difference of squares, sum/product form, completing the square}
%
% Arg 1 = the sheet's topic in one line. The website lists this instead of
%         "Sheet 03", so write the thing a reader would actually search for.
% Arg 2 = comma-separated, at most five: the tools this sheet introduces that
%         earlier sheets in the pillar did not. These become the website's tags.
%         Leave out anything the reader already met — the tags are meant to say
%         what is new today, not to summarise the sheet.
%
% Both are parsed straight out of this file by tools/build_website.py, so the
% sheet stays the single source of truth and the site cannot drift from it.
%
% This renders NOTHING. It is a declaration for the build script, not a piece of
% the sheet's layout: adding it to a sheet must not change that sheet's PDF, so
% the 25 existing sheets can be annotated without a single one of them being
% reissued. If the toolkit should also appear on the page, that is a separate
% decision about the sheet design — make it deliberately, here, not by leaving
% this macro printing something nobody asked it to.
\newcommand{\SpeedMeta}[2]{}

% ── Answer / Method / Investigate-further boxes (answer files only) ────────────
\newmdenv[
  backgroundcolor=lightgray,
  linecolor=medgray,
  linewidth=0.5pt,
  leftmargin=0pt, rightmargin=0pt,
  innerleftmargin=8pt, innerrightmargin=8pt,
  innertopmargin=5pt, innerbottommargin=5pt,
  skipabove=4pt, skipbelow=4pt
]{ansbox}

\newmdenv[
  backgroundcolor=white,
  linecolor=darkgray,
  linewidth=0.8pt,
  leftline=true, rightline=false, topline=false, bottomline=false,
  leftmargin=0pt, rightmargin=0pt,
  innerleftmargin=8pt, innerrightmargin=6pt,
  innertopmargin=4pt, innerbottommargin=4pt,
  skipabove=4pt, skipbelow=6pt
]{invbox}

\newcommand{\ans}[1]{%
  \begin{ansbox}
    \textbf{Answer:} #1
  \end{ansbox}}

\newcommand{\method}[1]{%
  {\small\textbf{Method:} #1}\par}

\newcommand{\inv}[1]{%
  \begin{invbox}
    \textbf{\small Investigate further:} {\small #1}
  \end{invbox}}

% ── Misc helpers ────────────────────────────────────────────────────────────────
\newcommand{\qn}[1]{\textbf{#1.}\;}

% Mistake-linking: attach a Top 5 Patterns / Common Traps bullet to the question
% label(s) in this sheet where it actually shows up, e.g. \seealso{B6, D2}
\newcommand{\seealso}[1]{\hfill\textit{\footnotesize(see #1)}}

% ── Graph options macro ─────────────────────────────────────────────────────────
% Usage: \GraphOptions{tikz code for A}{tikz code for B}{tikz code for C}{tikz code for D}
\newcommand{\GraphOptions}[4]{%
  \begin{center}
    \begin{tabular}{cc}
      \textbf{A)} & \textbf{B)} \\
      #1 & #2 \\[10pt]
      \textbf{C)} & \textbf{D)} \\
      #3 & #4
    \end{tabular}
  \end{center}
}

% ── Closing motivational rule + cross-promo footer (sheets only) ────────────────
% Usage: \SpeedClosing{"Quote text."}
\newcommand{\SpeedClosing}[1]{%
  \vspace{20pt}
  \begin{center}
  \rule{0.6\textwidth}{0.4pt}\\[4pt]
  {\small\itshape #1}\\[4pt]
  \rule{0.6\textwidth}{0.4pt}
  \end{center}
  \begin{center}
  {\small Found a mistake or want to contribute a sheet?}\\
  {\small Visit the open-source project at \href{https://github.com/The-CerealDev/speed-maths}{github.com/The-CerealDev/speed-maths}}
  \end{center}
}

% Editing THIS file changes the look of every sheet/answer across every pillar.
% ══════════════════════════════════════════════════════════════════════════════

\usepackage[a4paper, top=25mm, bottom=25mm, left=22mm, right=22mm]{geometry}
\usepackage{amsmath, amssymb}
\usepackage{enumitem}
\usepackage{booktabs}
\usepackage{array}
\usepackage{parskip}
\usepackage{microtype}
\usepackage{fancyhdr}
\usepackage{titlesec}
\usepackage{mdframed}
\usepackage{xcolor}
\usepackage[hidelinks]{hyperref}
\usepackage{graphicx}
\usepackage{tikz}
\usepackage{pgfplots}
\pgfplotsset{compat=1.18}

% ── Colours ───────────────────────────────────────────────────────────────────
\definecolor{darkgray}{gray}{0.15}
\definecolor{medgray}{gray}{0.45}
\definecolor{lightgray}{gray}{0.93}
\definecolor{boxgray}{gray}{0.88}

% ── Section formatting ──────────────────────────────────────────────────────────
\titleformat{\section}[block]
  {\normalfont\large\bfseries}{}
  {0pt}{}[\vspace{2pt}\hrule\vspace{6pt}]
\titlespacing*{\section}{0pt}{14pt}{4pt}

% ── List formatting ─────────────────────────────────────────────────────────────
\setlist[enumerate,1]{label=\textbf{\arabic*.}, leftmargin=2em, itemsep=10pt}

% ── Branding constants (edit here, not per-file) ────────────────────────────────
\newcommand{\SpeedExamLine}{\textbf{Speed Maths:} Competition \& Exam Prep}
\newcommand{\SpeedCredit}{Series by \href{https://www.linkedin.com/in/cerealdev/}{David Akinyele-Aje}}

% ── Header / footer ──────────────────────────────────────────────────────────────
% Usage (once, after % main (standalone preamble for editing)
% vim: nowrap: spell:
% ══════════════════════════════════════════════════════════════════════════════
% Speed Maths --- shared preamble
% Included by every sheet and answer file via:  \input{../../shared/preamble}
% Editing THIS file changes the look of every sheet/answer across every pillar.
% ══════════════════════════════════════════════════════════════════════════════

\usepackage[a4paper, top=25mm, bottom=25mm, left=22mm, right=22mm]{geometry}
\usepackage{amsmath, amssymb}
\usepackage{enumitem}
\usepackage{booktabs}
\usepackage{array}
\usepackage{parskip}
\usepackage{microtype}
\usepackage{fancyhdr}
\usepackage{titlesec}
\usepackage{mdframed}
\usepackage{xcolor}
\usepackage[hidelinks]{hyperref}
\usepackage{graphicx}
\usepackage{tikz}
\usepackage{pgfplots}
\pgfplotsset{compat=1.18}

% ── Colours ───────────────────────────────────────────────────────────────────
\definecolor{darkgray}{gray}{0.15}
\definecolor{medgray}{gray}{0.45}
\definecolor{lightgray}{gray}{0.93}
\definecolor{boxgray}{gray}{0.88}

% ── Section formatting ──────────────────────────────────────────────────────────
\titleformat{\section}[block]
  {\normalfont\large\bfseries}{}
  {0pt}{}[\vspace{2pt}\hrule\vspace{6pt}]
\titlespacing*{\section}{0pt}{14pt}{4pt}

% ── List formatting ─────────────────────────────────────────────────────────────
\setlist[enumerate,1]{label=\textbf{\arabic*.}, leftmargin=2em, itemsep=10pt}

% ── Branding constants (edit here, not per-file) ────────────────────────────────
\newcommand{\SpeedExamLine}{\textbf{Speed Maths:} Competition \& Exam Prep}
\newcommand{\SpeedCredit}{Series by \href{https://www.linkedin.com/in/cerealdev/}{David Akinyele-Aje}}

% ── Header / footer ──────────────────────────────────────────────────────────────
% Usage (once, after \input{../../shared/preamble} and before \begin{document}):
%   \SpeedHeader{Algebra}{3}
% First arg  = pillar name shown as "Speed <Pillar> --- Daily Drill"
% Second arg = worksheet number
\pagestyle{fancy}
\newcommand{\SpeedHeader}[2]{%
  \fancyhf{}%
  \renewcommand{\headrulewidth}{0.4pt}%
  \setlength{\headheight}{14pt}%
  \fancyhead[L]{\small\textbf{Speed #1 --- Daily Drill}}%
  \fancyhead[R]{\small Worksheet \##2}%
  \fancyfoot[L]{\small \SpeedExamLine}%
  \fancyfoot[C]{\small\thepage}%
  \fancyfoot[R]{\small \SpeedCredit}%
  \renewcommand{\footrulewidth}{0.4pt}%
}

% ── Title block (used at the top of every sheet AND answer file) ───────────────
% Usage: \SpeedTitleBlock{Daily Algebra Drill \#3}{Jane Doe}
% %% AGENTS: When generating a new sheet, ask the human contributor for the name they want credited in argument 2.
% For answer files, pass the "--- Answers \& Investigations" suffix in the arg.
\newcommand{\SpeedTitleBlock}[2]{%
  \begin{center}
    {\LARGE\bfseries #1}\\[4pt]
    {\normalsize\itshape Sheet by #2}\\[6pt]
    \hrule\vspace{4pt}
    \begin{tabular}{llcc}
      \toprule
      \textbf{Section} & \textbf{Topic} & \textbf{Questions} & \textbf{Target time}\\
      \midrule
      A & Rapid Recognition          & 10 & 2 min 30 sec \\
      B & Manipulation Drills        & 10 & 8 min        \\
      C & Substitution \& Structure  & 8  & 10 min       \\
      D & Challenge                  & 5  & 15 min       \\
      \bottomrule
    \end{tabular}
    \vspace{4pt}
    \hrule
  \end{center}
}

% ── Sheet metadata: topic + the day's new tools ────────────────────────────────
% Usage (immediately after \SpeedTitleBlock, sheets only):
%   \SpeedMeta{Expanding, factorising, and the identity toolkit}
%             {difference of squares, sum/product form, completing the square}
%
% Arg 1 = the sheet's topic in one line. The website lists this instead of
%         "Sheet 03", so write the thing a reader would actually search for.
% Arg 2 = comma-separated, at most five: the tools this sheet introduces that
%         earlier sheets in the pillar did not. These become the website's tags.
%         Leave out anything the reader already met — the tags are meant to say
%         what is new today, not to summarise the sheet.
%
% Both are parsed straight out of this file by tools/build_website.py, so the
% sheet stays the single source of truth and the site cannot drift from it.
%
% This renders NOTHING. It is a declaration for the build script, not a piece of
% the sheet's layout: adding it to a sheet must not change that sheet's PDF, so
% the 25 existing sheets can be annotated without a single one of them being
% reissued. If the toolkit should also appear on the page, that is a separate
% decision about the sheet design — make it deliberately, here, not by leaving
% this macro printing something nobody asked it to.
\newcommand{\SpeedMeta}[2]{}

% ── Answer / Method / Investigate-further boxes (answer files only) ────────────
\newmdenv[
  backgroundcolor=lightgray,
  linecolor=medgray,
  linewidth=0.5pt,
  leftmargin=0pt, rightmargin=0pt,
  innerleftmargin=8pt, innerrightmargin=8pt,
  innertopmargin=5pt, innerbottommargin=5pt,
  skipabove=4pt, skipbelow=4pt
]{ansbox}

\newmdenv[
  backgroundcolor=white,
  linecolor=darkgray,
  linewidth=0.8pt,
  leftline=true, rightline=false, topline=false, bottomline=false,
  leftmargin=0pt, rightmargin=0pt,
  innerleftmargin=8pt, innerrightmargin=6pt,
  innertopmargin=4pt, innerbottommargin=4pt,
  skipabove=4pt, skipbelow=6pt
]{invbox}

\newcommand{\ans}[1]{%
  \begin{ansbox}
    \textbf{Answer:} #1
  \end{ansbox}}

\newcommand{\method}[1]{%
  {\small\textbf{Method:} #1}\par}

\newcommand{\inv}[1]{%
  \begin{invbox}
    \textbf{\small Investigate further:} {\small #1}
  \end{invbox}}

% ── Misc helpers ────────────────────────────────────────────────────────────────
\newcommand{\qn}[1]{\textbf{#1.}\;}

% Mistake-linking: attach a Top 5 Patterns / Common Traps bullet to the question
% label(s) in this sheet where it actually shows up, e.g. \seealso{B6, D2}
\newcommand{\seealso}[1]{\hfill\textit{\footnotesize(see #1)}}

% ── Graph options macro ─────────────────────────────────────────────────────────
% Usage: \GraphOptions{tikz code for A}{tikz code for B}{tikz code for C}{tikz code for D}
\newcommand{\GraphOptions}[4]{%
  \begin{center}
    \begin{tabular}{cc}
      \textbf{A)} & \textbf{B)} \\
      #1 & #2 \\[10pt]
      \textbf{C)} & \textbf{D)} \\
      #3 & #4
    \end{tabular}
  \end{center}
}

% ── Closing motivational rule + cross-promo footer (sheets only) ────────────────
% Usage: \SpeedClosing{"Quote text."}
\newcommand{\SpeedClosing}[1]{%
  \vspace{20pt}
  \begin{center}
  \rule{0.6\textwidth}{0.4pt}\\[4pt]
  {\small\itshape #1}\\[4pt]
  \rule{0.6\textwidth}{0.4pt}
  \end{center}
  \begin{center}
  {\small Found a mistake or want to contribute a sheet?}\\
  {\small Visit the open-source project at \href{https://github.com/The-CerealDev/speed-maths}{github.com/The-CerealDev/speed-maths}}
  \end{center}
}
 and before \begin{document}):
%   \SpeedHeader{Algebra}{3}
% First arg  = pillar name shown as "Speed <Pillar> --- Daily Drill"
% Second arg = worksheet number
\pagestyle{fancy}
\newcommand{\SpeedHeader}[2]{%
  \fancyhf{}%
  \renewcommand{\headrulewidth}{0.4pt}%
  \setlength{\headheight}{14pt}%
  \fancyhead[L]{\small\textbf{Speed #1 --- Daily Drill}}%
  \fancyhead[R]{\small Worksheet \##2}%
  \fancyfoot[L]{\small \SpeedExamLine}%
  \fancyfoot[C]{\small\thepage}%
  \fancyfoot[R]{\small \SpeedCredit}%
  \renewcommand{\footrulewidth}{0.4pt}%
}

% ── Title block (used at the top of every sheet AND answer file) ───────────────
% Usage: \SpeedTitleBlock{Daily Algebra Drill \#3}{Jane Doe}
% %% AGENTS: When generating a new sheet, ask the human contributor for the name they want credited in argument 2.
% For answer files, pass the "--- Answers \& Investigations" suffix in the arg.
\newcommand{\SpeedTitleBlock}[2]{%
  \begin{center}
    {\LARGE\bfseries #1}\\[4pt]
    {\normalsize\itshape Sheet by #2}\\[6pt]
    \hrule\vspace{4pt}
    \begin{tabular}{llcc}
      \toprule
      \textbf{Section} & \textbf{Topic} & \textbf{Questions} & \textbf{Target time}\\
      \midrule
      A & Rapid Recognition          & 10 & 2 min 30 sec \\
      B & Manipulation Drills        & 10 & 8 min        \\
      C & Substitution \& Structure  & 8  & 10 min       \\
      D & Challenge                  & 5  & 15 min       \\
      \bottomrule
    \end{tabular}
    \vspace{4pt}
    \hrule
  \end{center}
}

% ── Sheet metadata: topic + the day's new tools ────────────────────────────────
% Usage (immediately after \SpeedTitleBlock, sheets only):
%   \SpeedMeta{Expanding, factorising, and the identity toolkit}
%             {difference of squares, sum/product form, completing the square}
%
% Arg 1 = the sheet's topic in one line. The website lists this instead of
%         "Sheet 03", so write the thing a reader would actually search for.
% Arg 2 = comma-separated, at most five: the tools this sheet introduces that
%         earlier sheets in the pillar did not. These become the website's tags.
%         Leave out anything the reader already met — the tags are meant to say
%         what is new today, not to summarise the sheet.
%
% Both are parsed straight out of this file by tools/build_website.py, so the
% sheet stays the single source of truth and the site cannot drift from it.
%
% This renders NOTHING. It is a declaration for the build script, not a piece of
% the sheet's layout: adding it to a sheet must not change that sheet's PDF, so
% the 25 existing sheets can be annotated without a single one of them being
% reissued. If the toolkit should also appear on the page, that is a separate
% decision about the sheet design — make it deliberately, here, not by leaving
% this macro printing something nobody asked it to.
\newcommand{\SpeedMeta}[2]{}

% ── Answer / Method / Investigate-further boxes (answer files only) ────────────
\newmdenv[
  backgroundcolor=lightgray,
  linecolor=medgray,
  linewidth=0.5pt,
  leftmargin=0pt, rightmargin=0pt,
  innerleftmargin=8pt, innerrightmargin=8pt,
  innertopmargin=5pt, innerbottommargin=5pt,
  skipabove=4pt, skipbelow=4pt
]{ansbox}

\newmdenv[
  backgroundcolor=white,
  linecolor=darkgray,
  linewidth=0.8pt,
  leftline=true, rightline=false, topline=false, bottomline=false,
  leftmargin=0pt, rightmargin=0pt,
  innerleftmargin=8pt, innerrightmargin=6pt,
  innertopmargin=4pt, innerbottommargin=4pt,
  skipabove=4pt, skipbelow=6pt
]{invbox}

\newcommand{\ans}[1]{%
  \begin{ansbox}
    \textbf{Answer:} #1
  \end{ansbox}}

\newcommand{\method}[1]{%
  {\small\textbf{Method:} #1}\par}

\newcommand{\inv}[1]{%
  \begin{invbox}
    \textbf{\small Investigate further:} {\small #1}
  \end{invbox}}

% ── Misc helpers ────────────────────────────────────────────────────────────────
\newcommand{\qn}[1]{\textbf{#1.}\;}

% Mistake-linking: attach a Top 5 Patterns / Common Traps bullet to the question
% label(s) in this sheet where it actually shows up, e.g. \seealso{B6, D2}
\newcommand{\seealso}[1]{\hfill\textit{\footnotesize(see #1)}}

% ── Graph options macro ─────────────────────────────────────────────────────────
% Usage: \GraphOptions{tikz code for A}{tikz code for B}{tikz code for C}{tikz code for D}
\newcommand{\GraphOptions}[4]{%
  \begin{center}
    \begin{tabular}{cc}
      \textbf{A)} & \textbf{B)} \\
      #1 & #2 \\[10pt]
      \textbf{C)} & \textbf{D)} \\
      #3 & #4
    \end{tabular}
  \end{center}
}

% ── Closing motivational rule + cross-promo footer (sheets only) ────────────────
% Usage: \SpeedClosing{"Quote text."}
\newcommand{\SpeedClosing}[1]{%
  \vspace{20pt}
  \begin{center}
  \rule{0.6\textwidth}{0.4pt}\\[4pt]
  {\small\itshape #1}\\[4pt]
  \rule{0.6\textwidth}{0.4pt}
  \end{center}
  \begin{center}
  {\small Found a mistake or want to contribute a sheet?}\\
  {\small Visit the open-source project at \href{https://github.com/The-CerealDev/speed-maths}{github.com/The-CerealDev/speed-maths}}
  \end{center}
}
 and before \begin{document}):
%   \SpeedHeader{Algebra}{3}
% First arg  = pillar name shown as "Speed <Pillar> --- Daily Drill"
% Second arg = worksheet number
\pagestyle{fancy}
\newcommand{\SpeedHeader}[2]{%
  \fancyhf{}%
  \renewcommand{\headrulewidth}{0.4pt}%
  \setlength{\headheight}{14pt}%
  \fancyhead[L]{\small\textbf{Speed #1 --- Daily Drill}}%
  \fancyhead[R]{\small Worksheet \##2}%
  \fancyfoot[L]{\small \SpeedExamLine}%
  \fancyfoot[C]{\small\thepage}%
  \fancyfoot[R]{\small \SpeedCredit}%
  \renewcommand{\footrulewidth}{0.4pt}%
}

% ── Title block (used at the top of every sheet AND answer file) ───────────────
% Usage: \SpeedTitleBlock{Daily Algebra Drill \#3}{Jane Doe}
% %% AGENTS: When generating a new sheet, ask the human contributor for the name they want credited in argument 2.
% For answer files, pass the "--- Answers \& Investigations" suffix in the arg.
\newcommand{\SpeedTitleBlock}[2]{%
  \begin{center}
    {\LARGE\bfseries #1}\\[4pt]
    {\normalsize\itshape Sheet by #2}\\[6pt]
    \hrule\vspace{4pt}
    \begin{tabular}{llcc}
      \toprule
      \textbf{Section} & \textbf{Topic} & \textbf{Questions} & \textbf{Target time}\\
      \midrule
      A & Rapid Recognition          & 10 & 2 min 30 sec \\
      B & Manipulation Drills        & 10 & 8 min        \\
      C & Substitution \& Structure  & 8  & 10 min       \\
      D & Challenge                  & 5  & 15 min       \\
      \bottomrule
    \end{tabular}
    \vspace{4pt}
    \hrule
  \end{center}
}

% ── Sheet metadata: topic + the day's new tools ────────────────────────────────
% Usage (immediately after \SpeedTitleBlock, sheets only):
%   \SpeedMeta{Expanding, factorising, and the identity toolkit}
%             {difference of squares, sum/product form, completing the square}
%
% Arg 1 = the sheet's topic in one line. The website lists this instead of
%         "Sheet 03", so write the thing a reader would actually search for.
% Arg 2 = comma-separated, at most five: the tools this sheet introduces that
%         earlier sheets in the pillar did not. These become the website's tags.
%         Leave out anything the reader already met — the tags are meant to say
%         what is new today, not to summarise the sheet.
%
% Both are parsed straight out of this file by tools/build_website.py, so the
% sheet stays the single source of truth and the site cannot drift from it.
%
% This renders NOTHING. It is a declaration for the build script, not a piece of
% the sheet's layout: adding it to a sheet must not change that sheet's PDF, so
% the 25 existing sheets can be annotated without a single one of them being
% reissued. If the toolkit should also appear on the page, that is a separate
% decision about the sheet design — make it deliberately, here, not by leaving
% this macro printing something nobody asked it to.
\newcommand{\SpeedMeta}[2]{}

% ── Answer / Method / Investigate-further boxes (answer files only) ────────────
\newmdenv[
  backgroundcolor=lightgray,
  linecolor=medgray,
  linewidth=0.5pt,
  leftmargin=0pt, rightmargin=0pt,
  innerleftmargin=8pt, innerrightmargin=8pt,
  innertopmargin=5pt, innerbottommargin=5pt,
  skipabove=4pt, skipbelow=4pt
]{ansbox}

\newmdenv[
  backgroundcolor=white,
  linecolor=darkgray,
  linewidth=0.8pt,
  leftline=true, rightline=false, topline=false, bottomline=false,
  leftmargin=0pt, rightmargin=0pt,
  innerleftmargin=8pt, innerrightmargin=6pt,
  innertopmargin=4pt, innerbottommargin=4pt,
  skipabove=4pt, skipbelow=6pt
]{invbox}

\newcommand{\ans}[1]{%
  \begin{ansbox}
    \textbf{Answer:} #1
  \end{ansbox}}

\newcommand{\method}[1]{%
  {\small\textbf{Method:} #1}\par}

\newcommand{\inv}[1]{%
  \begin{invbox}
    \textbf{\small Investigate further:} {\small #1}
  \end{invbox}}

% ── Misc helpers ────────────────────────────────────────────────────────────────
\newcommand{\qn}[1]{\textbf{#1.}\;}

% Mistake-linking: attach a Top 5 Patterns / Common Traps bullet to the question
% label(s) in this sheet where it actually shows up, e.g. \seealso{B6, D2}
\newcommand{\seealso}[1]{\hfill\textit{\footnotesize(see #1)}}

% ── Graph options macro ─────────────────────────────────────────────────────────
% Usage: \GraphOptions{tikz code for A}{tikz code for B}{tikz code for C}{tikz code for D}
\newcommand{\GraphOptions}[4]{%
  \begin{center}
    \begin{tabular}{cc}
      \textbf{A)} & \textbf{B)} \\
      #1 & #2 \\[10pt]
      \textbf{C)} & \textbf{D)} \\
      #3 & #4
    \end{tabular}
  \end{center}
}

% ── Closing motivational rule + cross-promo footer (sheets only) ────────────────
% Usage: \SpeedClosing{"Quote text."}
\newcommand{\SpeedClosing}[1]{%
  \vspace{20pt}
  \begin{center}
  \rule{0.6\textwidth}{0.4pt}\\[4pt]
  {\small\itshape #1}\\[4pt]
  \rule{0.6\textwidth}{0.4pt}
  \end{center}
  \begin{center}
  {\small Found a mistake or want to contribute a sheet?}\\
  {\small Visit the open-source project at \href{https://github.com/The-CerealDev/speed-maths}{github.com/The-CerealDev/speed-maths}}
  \end{center}
}

\SpeedHeader{Combinatorics}{4}

\begin{document}

\SpeedTitleBlock{Daily Combinatorics Drill \#4 --- Answers \& Investigations}
\noindent\textit{New toolkit today: Binomial Theorem $\cdot$ Pascal's Rule $\cdot$ Coefficient Extraction $\cdot$ Substituting Values into Expansions.}

\section*{Section A \quad Rapid Recognition \hfill{\normalfont\small($\leq$15\,s each)}}
%═══════════════════════════════════════════════════════════════════════════════

\begin{enumerate}

%── A1 ──────────────────────────────────────────────────────────────────────────
\item Coefficient of $x^2$ in $(1+x)^5$.

\ans{$\dbinom{5}{2}=10$}
\method{Binomial theorem: the $x^k$ coefficient in $(1+x)^n$ is $\binom{n}{k}$.}
\inv{\emph{Why} is it $\binom{5}{2}$? Expanding $(1+x)(1+x)(1+x)(1+x)(1+x)$, an $x^2$ term arises by choosing the $x$ from exactly 2 of the 5 brackets. The binomial theorem \emph{is} a counting statement --- that lens powers everything on this sheet.}

%── A2 ──────────────────────────────────────────────────────────────────────────
\item Coefficient of $x$ in $(x+2)^3$.

\ans{$12$}
\method{Term with $x^1$: $\binom{3}{1}x\cdot2^2=12x$. Keep track of which quantity is being raised to which power.}
\inv{Write out all four terms of $(x+2)^3$ and check they sum to $27$ at $x=1$. The ``evaluate at $x=1$'' check catches most expansion slips in exams.}

%── A3 ──────────────────────────────────────────────────────────────────────────
\item Row 5 of Pascal's triangle.

\ans{$1,\;5,\;10,\;10,\;5,\;1$}
\method{Adjacent sums of row 4 ($1,4,6,4,1$), or $\binom{5}{k}$ directly.}
\inv{Memorise rows 0--8 cold. Then look at the triangle mod 2 (odd entries only) --- shade them for rows 0--16 and meet the Sierpi\'nski triangle. D3 explains the pattern.}

%── A4 ──────────────────────────────────────────────────────────────────────────
\item Coefficient of $x^3y^2$ in $(x+y)^5$.

\ans{$\dbinom{5}{2}=10$}
\method{Choose which 2 of the 5 brackets contribute $y$: $\binom{5}{2}$.}
\inv{Coefficients of $x^3y^2$ and $x^2y^3$ are equal --- symmetry of $\binom{5}{2}=\binom{5}{3}$ made visible. In how many terms does $(x+y)^5$ expand? (6 --- one per power split.)}

%── A5 ──────────────────────────────────────────────────────────────────────────
\item Sum of all coefficients of $(1+x)^7$.

\ans{$2^7=128$}
\method{Substitute $x=1$: the expansion collapses to its coefficient sum.}
\inv{Substituting values into an identity is the sheet's master move. What does $x=2$ give? ($3^7$ --- the sum $\sum\binom{7}{k}2^k$.) Every substitution is a new identity for free.}

%── A6 ──────────────────────────────────────────────────────────────────────────
\item Constant term of $\left(x+\frac{1}{x}\right)^4$.

\ans{$\dbinom{4}{2}=6$}
\method{Powers cancel when the $x$'s and $\frac1x$'s balance: 2 of each.}
\inv{Why does $\left(x+\frac1x\right)^5$ have \emph{no} constant term? (Odd total degree --- balance is impossible.) Parity kills entire terms before any computation.}

%── A7 ──────────────────────────────────────────────────────────────────────────
\item Express $\dbinom{9}{4}+\dbinom{9}{5}$ as one coefficient.

\ans{$\dbinom{10}{5}$\quad(which is 252)}
\method{Pascal's rule: adjacent entries sum to the entry below. No evaluation needed --- the recognition is the answer.}
\inv{Prove it in one line with committees (a committee of 5 from 10 either contains the distinguished person or not) --- you will need exactly this argument in C8.}

%── A8 ──────────────────────────────────────────────────────────────────────────
\item $\dbinom{6}{0}-\dbinom{6}{1}+\dbinom{6}{2}-\cdots+\dbinom{6}{6}$.

\ans{$0$}
\method{Substitute $x=-1$ into $(1+x)^6$: $0^6=0$.}
\inv{Consequence: in any row, the even-index entries and odd-index entries have equal sums (each $2^{n-1}$ --- B5 uses this). Pair with sheet 2's subset-parity bijection: same fact, two proofs.}

%── A9 ──────────────────────────────────────────────────────────────────────────
\item True or false: coefficient of $x^2$ in $(1-x)^{10}$ is negative.

\ans{False --- it is $+45$}
\method{The $x^2$ term is $\binom{10}{2}(-x)^2=45x^2$: the square eats the sign. Signs alternate, and even powers are positive.}
\inv{TMUA-style sign traps live exactly here. Which coefficients of $(1-x)^{10}$ \emph{are} negative? (The odd ones.) State the general sign rule for $(a-b)^n$ and never lose a mark to it again.}

%── A10 ─────────────────────────────────────────────────────────────────────────
\item Largest coefficient in $(1+x)^4$.

\ans{$\dbinom{4}{2}=6$}
\method{Rows of Pascal's triangle rise to the middle and fall symmetrically.}
\inv{For even $n$ there is one maximum; for odd $n$, two equal ones. Prove the rise-then-fall by examining the ratio $\binom{n}{k+1}/\binom{n}{k}=\frac{n-k}{k+1}$ --- when does it drop below 1? (C2 uses this.)}

\end{enumerate}

%═══════════════════════════════════════════════════════════════════════════════
\section*{Section B \quad Manipulation Drills \hfill{\normalfont\small($\leq$50\,s each)}}
%═══════════════════════════════════════════════════════════════════════════════

\begin{enumerate}

%── B1 ──────────────────────────────────────────────────────────────────────────
\item Coefficient of $x^3$ in $(2+3x)^5$.

\ans{$\dbinom{5}{3}2^2 3^3=10\cdot4\cdot27=1\,080$}
\method{General term $\binom{5}{k}2^{5-k}(3x)^k$; set $k=3$. Both constants carry powers --- forgetting $2^{5-k}$ is the classic slip.}
\inv{Check digit-cheaply: at $x=1$ the whole expansion is $5^5=3125$; your five coefficients must sum to it. Run the check --- it takes 20 seconds and catches the missing-power error instantly.}

%── B2 ──────────────────────────────────────────────────────────────────────────
\item Coefficient of $x^5$ in $(1+2x)^4(1+x)^3$.

\ans{$168$}
\method{Convolve the tails only: from $(1+2x)^4$ take coefficients $16, 32, 24$ (of $x^4,x^3,x^2$) and from $(1+x)^3$ take $3, 3, 1$ (of $x,x^2,x^3$): $16\cdot3+32\cdot3+24\cdot1=168$.}
\inv{Had both brackets shared the base $(1+x)$, you would merge first: $(1+x)^4(1+x)^3=(1+x)^7$, coefficient $\binom{7}{5}=21$ in five seconds. Check the bases before you convolve --- and note that separate-then-convolve on the shared-base version proves $\sum_k\binom{4}{k}\binom{3}{5-k}=\binom{7}{5}$, a Vandermonde identity (D2).}

%── B3 ──────────────────────────────────────────────────────────────────────────
\item Term independent of $x$ in $\left(2x+\frac{1}{x^2}\right)^6$.

\ans{$\dbinom{6}{2}2^4=240$}
\method{General term: $\binom{6}{k}(2x)^{6-k}x^{-2k}$ has degree $6-3k$. Zero degree forces $k=2$: $\binom{6}{2}2^4$.}
\inv{Degree bookkeeping first, arithmetic second. Which powers of $x$ \emph{can} occur here? ($6-3k$: only $6,3,0,-3,-6,-9,-12$ --- gaps of 3.) Knowing the degree lattice tells you instantly that e.g.\ an $x^2$ term is impossible.}

%── B4 ──────────────────────────────────────────────────────────────────────────
\item Coefficient of $x^2$ in $(1+x)^n$ is 66. Find $n$.

\ans{$n=12$}
\method{$\binom{n}{2}=66\Rightarrow n(n-1)=132=12\cdot11$.}
\inv{Sheet 2's consecutive-product reflex again. Harder inverse: the coefficient of $x^3$ is 220 --- find $n$ ($n(n-1)(n-2)=1320=12\cdot11\cdot10$). Estimation ($\sqrt[3]{1320}\approx11$) then check.}

%── B5 ──────────────────────────────────────────────────────────────────────────
\item Coefficient of $x^6$ in $\left(x^2+\frac{1}{x}\right)^6$.

\ans{$\dbinom{6}{2}=15$}
\method{General term $\binom{6}{k}(x^2)^{6-k}x^{-k}$ has degree $12-3k$; degree 6 forces $k=2$.}
\inv{The degree lattice here is $12,9,6,3,0,-3,-6$. Which is the constant term? ($k=4$: $\binom{6}{4}=15$ --- equal to the $x^6$ coefficient. Coincidence, or a symmetry of this particular bracket? Substitute $x\to\frac{1}{x}$ and find out.)}

%── B6 ──────────────────────────────────────────────────────────────────────────
\item Coefficient of $x^2$ in $(1+2x)^4(1-x)^2$.

\ans{$9$}
\method{Convolve only what is needed: $x^2$ coefficient $=24\cdot1+8\cdot(-2)+1\cdot1=9$, using $(1+2x)^4=1+8x+24x^2+\cdots$ and $(1-x)^2=1-2x+x^2$.}
\inv{Never expand whole products for one coefficient --- harvest each factor only up to the degree you need. Practise: coefficient of $x^3$ in the same product ($32\cdot1+24\cdot(-2)+8\cdot1=-8$).}

%── B7 ──────────────────────────────────────────────────────────────────────────
\item Evaluate $11^3$ and $9^3$ by binomial.

\ans{$11^3=1\,331$;\quad $9^3=729$}
\method{$(10+1)^3=1000+300+30+1$; $(10-1)^3=1000-300+30-1$. Pascal row 3 does the arithmetic.}
\inv{Notice $11^3$'s digits ARE row 3 of Pascal's triangle (1,3,3,1) --- and $11^4=14641$ is row 4. When does the pattern break, and why? (Row 5 has a two-digit entry: carrying destroys the display.)}

%── B8 ──────────────────────────────────────────────────────────────────────────
\item Coefficient of $x^3$ in $(1+kx)^6$ is 160. Find $k$.

\ans{$k=2$}
\method{$\binom{6}{3}k^3=20k^3=160\Rightarrow k^3=8$.}
\inv{Could a negative or fractional $k$ ever be the intended answer in such questions? Solve $\binom{6}{3}k^3=-540$ and $\binom{4}{2}k^2=\frac{3}{2}$ to see both arise naturally in TMUA-style settings.}

%── B9 ──────────────────────────────────────────────────────────────────────────
\item $\dbinom{n}{2}=\dbinom{n}{3}$. Find $n$.

\ans{$n=5$}
\method{Equal entries sit symmetrically: $\binom{n}{r}=\binom{n}{s}$ (with $r\neq s$) forces $r+s=n$. So $n=2+3=5$.}
\inv{Solve without symmetry, by ratio: $\frac{\binom{n}{3}}{\binom{n}{2}}=\frac{n-2}{3}=1$. The ratio form generalises to ``for which $k$ is $\binom{n}{k}$ maximal?'' --- close the loop with A10.}

%── B10 ─────────────────────────────────────────────────────────────────────────
\item Coefficient of $a^2b^2c$ in $(a+b+c)^5$. \textit{\small(cf.\ TMUA 2020 P1 Q13)}

\ans{$\dfrac{5!}{2!\,2!\,1!}=30$}
\method{Multinomial: distribute the 5 brackets among the letters --- exactly the BANANA formula from sheet 3 wearing algebra.}
\inv{Repeated-letter words, flag rows, multinomial coefficients: one formula, three sheets. How many \emph{distinct monomials} does $(a+b+c)^5$ have? (Solutions of $i+j+k=5$: that is sheet 5's stars-and-bars, arriving tomorrow.)}

\end{enumerate}

%═══════════════════════════════════════════════════════════════════════════════
\section*{Section C \quad Substitution \& Structure \hfill{\normalfont\small($\leq$90\,s each)}}
%═══════════════════════════════════════════════════════════════════════════════

\begin{enumerate}

%── C1 ──────────────────────────────────────────────────────────────────────────
\item $\dbinom{10}{0}+\dbinom{10}{2}+\cdots+\dbinom{10}{10}$.

\ans{$2^9=512$}
\method{Substitute twice into $(1+x)^{10}$: $x=1$ gives $2^{10}$ (all terms), $x=-1$ gives $0$ (alternating). Adding kills the odd-index terms and doubles the even ones: sum $=(2^{10}+0)/2=512$.}
\inv{The add/subtract-substitutions trick extracts any alternating pattern. Try extracting $\binom{10}{0}+\binom{10}{4}+\binom{10}{8}$ --- what would you substitute? (Fourth roots of unity: a genuine BMO-level extension to sit with.)}

%── C2 ──────────────────────────────────────────────────────────────────────────
\item Greatest coefficient in $(1+x)^8$.

\ans{$\dbinom{8}{4}=70$}
\method{Ratio test from A10's inv: $\frac{\binom{8}{k+1}}{\binom{8}{k}}=\frac{8-k}{k+1}>1$ exactly while $k<3.5$, so coefficients climb to $k=4$ then fall.}
\inv{Now find the greatest coefficient of $(1+2x)^8$ --- the ratio becomes $\frac{2(8-k)}{k+1}$ and the peak moves right (to $k=5$: check it). Peaks shift towards the larger term; TMUA loves this.}

%── C3 ──────────────────────────────────────────────────────────────────────────
\item Estimate $1.01^{10}$ from three terms.

\ans{$1+10(0.01)+45(0.01)^2=1.1045$}
\method{$(1+x)^{10}$ at $x=0.01$; terms shrink so fast the tail is invisible at 4 d.p.}
\inv{Bound the error: the next term is $120(0.01)^3=1.2\times10^{-4}$... which \emph{would} touch the 4th decimal place. Investigate: the true value is $1.10462\ldots$, so why did the estimate still round correctly? (The tail is positive but small enough. Always \emph{bound}, not assume.)}

%── C4 ──────────────────────────────────────────────────────────────────────────
\item In $(1+3x)^n$, coefficient of $x^5$ is three times that of $x^4$. Find $n$.

\ans{$n=9$}
\method{$\binom{n}{5}3^5=3\cdot\binom{n}{4}3^4$. The powers of 3 cancel completely ($3^5$ on both sides), leaving $\binom{n}{5}=\binom{n}{4}$, so $n=4+5=9$ by symmetry.}
\inv{The constants were a decoy --- strip them before solving. Equal \emph{adjacent} plain coefficients happen only for odd $n$ at the centre (A10's two equal maxima). For a genuine constant effect, solve: in $(1+2x)^n$ the $x^5$ coefficient is \emph{four} times the $x^4$ one; now the 2's don't cancel.}

%── C5 ──────────────────────────────────────────────────────────────────────────
\item Coefficient of $x^4$ in $(1+x)^5(1-2x)^3$. \textit{\small(cf.\ MAT 2023 Q1F)}

\ans{$25$}
\method{Convolution with signs, harvesting only degree-4 pairs: $5\cdot1+10\cdot(-6)+10\cdot12+5\cdot(-8)=5-60+120-40=25$. Lay the two coefficient lists side by side and zip them --- do not expand the product.}
\inv{MAT's Q1 multiple-choice regularly hides one sign or one binomial power slip among its options. Recompute with the factors swapped in your zip to double-check --- convolution is symmetric, so any asymmetry in your two answers locates the error.}

%── C6 ──────────────────────────────────────────────────────────────────────────
\item $\displaystyle\sum_{k=0}^{5}k\binom{5}{k}$.

\ans{$5\cdot2^4=80$}
\method{Committee--chair (sheet 2, B5): $k\binom{5}{k}=5\binom{4}{k-1}$, and summing the right side gives $5\cdot2^4$. Alternatively differentiate $(1+x)^5$ and set $x=1$.}
\inv{Two proofs from two different sheets --- the counting one and the calculus one. Use either to evaluate $\sum k^2\binom{5}{k}$ (answer $5\cdot6\cdot2^3=$ check it: $k^2=k(k-1)+k$). Identities stack.}

%── C7 ──────────────────────────────────────────────────────────────────────────
\item Evaluate $\displaystyle\sum_{k=0}^{6}\binom{6}{k}2^k$ by substitution.

\ans{$3^6=729$}
\method{The sum is exactly $(1+x)^6$ at $x=2$: $(1+2)^6=729$. Recognising a binomial expansion \emph{in reverse} is the skill --- the $2^k$ pattern names the substitution.}
\inv{Every value of $x$ turns $(1+x)^n$ into an identity: $x=10$ explains why row-$n$ digits ``read as'' $11^n$ (B7), and $x=-\frac12$ gives an alternating shrinking sum. Evaluate $\sum\binom{6}{k}(-1)^k2^{6-k}$ by spotting \emph{its} disguise ($(2-1)^6=1$).}

%── C8 ──────────────────────────────────────────────────────────────────────────
\item Prove Pascal's rule by committee; build row 7 from row 6.

\ans{Row 7: $1,7,21,35,35,21,7,1$}
\method{\emph{Proof.} A committee of $r$ from $n$ people either includes a distinguished person $P$ or not. Including $P$: $\binom{n-1}{r-1}$ choices for the rest. Excluding $P$: $\binom{n-1}{r}$. These cases are disjoint and exhaustive, so $\binom{n}{r}=\binom{n-1}{r-1}+\binom{n-1}{r}$. $\square$ Adjacent sums of row 6 ($1,6,15,20,15,6,1$) then give row 7.}
\inv{This one-sentence proof is a model of BMO1 writing: cases named, disjointness noted, conclusion drawn. Rewrite sheet 2 B9's Priya split as a formal proof in the same register --- Section D now expects this standard.}

\end{enumerate}

%═══════════════════════════════════════════════════════════════════════════════
\section*{Section D \quad Challenge \hfill{\normalfont\small(SMC / BMO1 difficulty)}}
%═══════════════════════════════════════════════════════════════════════════════

\begin{enumerate}

%── D1 ──────────────────────────────────────────────────────────────────────────
\item Prove $\dbinom{2n}{n}$ is even for all $n\geq1$.

\ans{Proof below.}
\method{\emph{Proof.} By Pascal's rule, $\binom{2n}{n}=\binom{2n-1}{n-1}+\binom{2n-1}{n}$. By the symmetry $\binom{m}{r}=\binom{m}{m-r}$ applied with $m=2n-1$: $\binom{2n-1}{n-1}=\binom{2n-1}{(2n-1)-(n-1)}=\binom{2n-1}{n}$. So $\binom{2n}{n}$ is the sum of two \emph{equal} integers, hence even. $\square$}
\method{\emph{Alternative (bijective) proof.} Pair each $n$-subset $S$ of $\{1,\ldots,2n\}$ with its complement $S^c$. Since $|S|=n=|S^c|$ and $S\neq S^c$ always, the $n$-subsets split into disjoint pairs, so their number is even. $\square$}
\inv{Two complete written proofs --- study how each names its key step. Sharper question: what is the largest power of 2 dividing $\binom{2n}{n}$? Compute it for $n=1,\ldots,8$ and conjecture (Kummer's theorem: the number of carries when adding $n+n$ in binary).}

%── D2 ──────────────────────────────────────────────────────────────────────────
\item Show $\displaystyle\sum_{k=0}^{10}\binom{10}{k}^2=\binom{20}{10}$.

\ans{Both count $\binom{20}{10}=184\,756$.}
\method{\emph{Proof.} Consider choosing 10 people from a room of 10 women and 10 men: $\binom{20}{10}$ ways. Alternatively, classify each choice by the number $k$ of women chosen: $\binom{10}{k}$ ways for the women and $\binom{10}{10-k}$ for the men, and $\binom{10}{10-k}=\binom{10}{k}$ by symmetry. Summing the disjoint classes: $\sum_k\binom{10}{k}^2$. Two counts of one set are equal. $\square$}
\inv{This is Vandermonde's identity $\sum_k\binom{m}{k}\binom{n}{r-k}=\binom{m+n}{r}$ at $m=n=r=10$ --- sheet 2's C5 inv, now proved properly. Also prove it by extracting the $x^{10}$ coefficient from $(1+x)^{10}(1+x)^{10}$: the algebra proof and the committee proof should feel like translations of each other.}

%── D3 ──────────────────────────────────────────────────────────────────────────
\item Prove exactly four of the $\dbinom{12}{k}$ are odd. \textit{\small(cf.\ TMUA 2023 P1 Q6)}

\ans{The odd ones are exactly $k=0,4,8,12$.}
\method{\emph{Proof.} All congruences are mod 2. First, $(1+x)^2=1+2x+x^2\equiv1+x^2$, and squaring again, $(1+x)^4\equiv(1+x^2)^2\equiv1+x^4$. Hence
$(1+x)^{12}=\left((1+x)^4\right)^3\equiv(1+x^4)^3=1+3x^4+3x^8+x^{12}\equiv1+x^4+x^8+x^{12}.$
Comparing coefficients of $x^k$ on both sides: $\binom{12}{k}$ is odd precisely when $k\in\{0,4,8,12\}$ --- four values. $\square$}
\inv{General law (Lucas): $\binom{n}{k}$ is odd iff every binary digit of $k$ is $\leq$ the corresponding digit of $n$. Here $n=12=1100_2$, so $k\in\{0000,0100,1000,1100\}$ --- the same four. Verify Lucas on row 10 ($1010_2$: which $k$?), and connect to A3's Sierpi\'nski picture.}

%── D4 ──────────────────────────────────────────────────────────────────────────
\item Show $\dbinom{2}{2}+\dbinom{3}{2}+\cdots+\dbinom{9}{2}=\dbinom{10}{3}$.

\ans{$\dbinom{10}{3}=120$}
\method{\emph{Proof (classify by largest element).} Count the 3-subsets of $\{1,2,\ldots,10\}$: there are $\binom{10}{3}$. Classify each subset by its largest element $m+1$ (so $m+1$ ranges over $3,\ldots,10$): the other two elements are then any 2-subset of $\{1,\ldots,m\}$, giving $\binom{m}{2}$ subsets in that class. The classes are disjoint and exhaustive, so $\sum_{m=2}^{9}\binom{m}{2}=\binom{10}{3}$. $\square$}
\inv{This is the \emph{hockey stick identity}: any diagonal run in Pascal's triangle sums to the entry below-right of its end. Prove it a second way by iterating Pascal's rule. Then use it to re-derive $1+2+\cdots+n=\binom{n+1}{2}$ and $\sum m^2$ (write $m^2=2\binom{m}{2}+\binom{m}{1}$).}

%── D5 ──────────────────────────────────────────────────────────────────────────
\item $p$ prime, $0<k<p$: prove $p\mid\dbinom{p}{k}$.

\ans{Proof below.}
\method{\emph{Proof.} Write $k!\,\binom{p}{k}=p(p-1)\cdots(p-k+1)$. The right side is visibly divisible by $p$, so $p\mid k!\binom{p}{k}$. Since $p$ is prime and every factor of $k!=1\cdot2\cdots k$ is less than $p$, we have $p\nmid k!$, so $p$ and $k!$ are coprime. A prime dividing a product must divide one of the factors; as it cannot divide $k!$, it divides $\binom{p}{k}$. $\square$ (Note where primality is used --- \emph{twice}.)}
\inv{Consequence: $(1+x)^p\equiv1+x^p\pmod p$ --- the ``freshman's dream'' is a theorem mod $p$. Feed $x=1$ into it to get $2^p\equiv2\pmod p$: you have proved Fermat's little theorem for base 2. Test the failure for composites: is $\binom{4}{2}$ divisible by 4? This lemma returns as a workhorse in the number-theory pillar.}

\end{enumerate}

%═══════════════════════════════════════════════════════════════════════════════
\section*{Top 5 Patterns Today}
%═══════════════════════════════════════════════════════════════════════════════

\begin{enumerate}[leftmargin=2em, itemsep=4pt]
  \item \textbf{The general term.} Write $\binom{n}{k}a^{n-k}b^k$ \emph{with the constants inside}, solve the degree equation for $k$, then compute. Never expand what you can index.
  \item \textbf{Substitute values into identities.} $x=1$: coefficient sum. $x=-1$: alternating sum. Add/subtract the two: even/odd extraction. Every substitution is a theorem.
  \item \textbf{Pascal's rule = in-or-out.} $\binom{n}{r}=\binom{n-1}{r-1}+\binom{n-1}{r}$ is the committee split from sheet 2, and it generates the whole triangle.
  \item \textbf{Count one thing two ways.} Vandermonde, hockey stick, committee--chair: pick the right scenario and the identity proves itself in a sentence.
  \item \textbf{Binomials mod small primes.} $(1+x)^2\equiv1+x^2\pmod2$ propagates to $(1+x)^{2^m}\equiv1+x^{2^m}$: parity patterns of whole rows fall out (Lucas/Sierpi\'nski).
\end{enumerate}

%═══════════════════════════════════════════════════════════════════════════════
\section*{Common Traps to Avoid}
%═══════════════════════════════════════════════════════════════════════════════

\begin{enumerate}[leftmargin=2em, itemsep=4pt]
  \item \textbf{Dropping the constant's power.} In $(2+3x)^5$ the $2^{5-k}$ is compulsory. Coefficient-sum check at $x=1$ catches this in seconds.
  \item \textbf{Sign slips with $(a-b)^n$.} Keep $(-b)^k$ intact inside the general term; decide the sign from the parity of $k$, never by feel.
  \item \textbf{Expanding both factors fully for one coefficient.} Harvest each factor only to the degree needed and zip. Full expansion is slow and error-prone.
  \item \textbf{Assuming symmetric-looking sums need algebra.} $\sum\binom{10}{k}^2$ falls to a story about choosing 10 people. Reach for the committee before the formula sheet.
  \item \textbf{Using primality nowhere.} If your proof of D5 never invokes that $p$ is prime, it proves the false statement $4\mid\binom{4}{2}$. Locate where every hypothesis earns its keep.
\end{enumerate}

\SpeedClosing{``An identity is a story told twice ---
once for each side of the equals sign.''}

\end{document}
